\section*{Bab \ref{ch:g11:redox} --- Oksidasi dan Reduksi}

\begin{solution}{exo:g11:redox:1}
\omterm{def:g11:redox:oxidant}{Oksidator}: spesies yang mampu menerima \omterm{def:g9:inside-the-atom:nucleus}{elektron}. \omterm{def:g11:redox:oxidant}{Reduktor}: spesies yang
mampu memberikan \omterm{def:g9:inside-the-atom:nucleus}{elektron}. \omterm{def:g11:redox:oxidation}{Oksidasi}: pelepasan \omterm{def:g9:inside-the-atom:nucleus}{elektron}. \omterm{def:g11:redox:oxidation}{Reduksi}:
penerimaan \omterm{def:g9:inside-the-atom:nucleus}{elektron}.
\end{solution}

\begin{solution}{exo:g11:redox:2}
\ce{Fe^{2+}}/\ce{Fe}; \ce{Ag+}/\ce{Ag}; \ce{I2}/\ce{I-};
\ce{Fe^{3+}}/\ce{Fe^{2+}}.
\end{solution}

\begin{solution}{exo:g11:redox:3}
\ce{Ag+ + e- <=> Ag}; \ce{Al^{3+} + 3e- <=> Al};
\ce{Fe^{3+} + e- <=> Fe^{2+}}; \ce{Cl2 + 2e- <=> 2Cl-}.
\end{solution}

\begin{solution}{exo:g11:redox:4}
\omterm{def:g11:redox:oxidation}{Oksidasi} (\omterm{def:g9:inside-the-atom:nucleus}{elektron} dilepaskan); besi bertindak sebagai \omterm{def:g11:redox:oxidant}{reduktor}.
\end{solution}

\begin{solution}{exo:g11:redox:5}
Seng dioksidasi dan merupakan \omterm{def:g11:redox:oxidant}{reduktor}; \ce{Cu^{2+}} direduksi dan
merupakan \omterm{def:g11:redox:oxidant}{oksidator}. Dua \omterm{def:g9:inside-the-atom:nucleus}{elektron} per \omterm{def:g7:atoms-and-molecules:atom}{atom} seng.
\end{solution}

\begin{solution}{exo:g11:redox:6}
Mulai dari \ce{H2O2 -> H2O}: % ce-unbalanced-ok (step 1 of the method)
oksigen: dua O di kiri, jadi \ce{2H2O} di kanan; hidrogen: 2 di kiri,
4 di kanan, jadi \ce{2H+} di kiri; muatan: $+2$ di kiri, 0 di kanan,
jadi dua \omterm{def:g9:inside-the-atom:nucleus}{elektron} di kiri:
\ce{H2O2 + 2H+ + 2e- <=> 2H2O}.
\end{solution}

\begin{solution}{exo:g11:redox:7}
\ce{I2 + 2e- -> 2I-} dan \ce{2S2O3^{2-} -> S4O6^{2-} + 2e-}; masing-masing
dua \omterm{def:g9:inside-the-atom:nucleus}{elektron}, jadi \ce{I2 + 2S2O3^{2-} -> 2I- + S4O6^{2-}}.
\end{solution}

\begin{solution}{exo:g11:redox:8}
\ce{Cu -> Cu^{2+} + 2e-} dan $2\times$ (\ce{Ag+ + e- -> Ag}):
\ce{Cu + 2Ag+ -> Cu^{2+} + 2Ag}. \omterm{def:g9:ions:ion}{Ion-ion} \ce{Cu^{2+}} yang terbentuk
berwarna biru.
\end{solution}

\begin{solution}{exo:g11:redox:9}
\ce{C6H8O6 -> C6H6O6 + 2H+ + 2e-} dan \ce{I2 + 2e- -> 2I-}:
\ce{C6H8O6 + I2 -> C6H6O6 + 2I- + 2H+}. Vitamin C adalah \omterm{def:g11:redox:oxidant}{reduktornya}.
\end{solution}

\begin{solution}{exo:g11:redox:10}
\ce{Al -> Al^{3+} + 3e-} memberikan 3 \omterm{def:g9:inside-the-atom:nucleus}{elektron}, \ce{Cu^{2+} + 2e- -> Cu}
mengambil 2; bilangan persekutuan terkecilnya 6, jadi yang pertama
dikalikan 2 dan yang kedua 3: \ce{2Al + 3Cu^{2+} -> 2Al^{3+} + 3Cu}.
\end{solution}

\begin{solution}{exo:g11:redox:11}
\ce{H2O2 + 2H+ + 2e- -> 2H2O} dan $2\times$ (\ce{Fe^{2+} -> Fe^{3+} + e-}):
\ce{H2O2 + 2H+ + 2Fe^{2+} -> 2H2O + 2Fe^{3+}}.
\end{solution}

\begin{solution}{exo:g11:redox:12}
$n(\ce{Zn}) = 1.30/65.4 = \qty{1.99e-2}{mol}$. Satu \omterm{def:g7:atoms-and-molecules:atom}{atom} tembaga
mengendap untuk setiap \omterm{def:g7:atoms-and-molecules:atom}{atom} seng yang dioksidasi, jadi
$n(\ce{Cu}) = \qty{1.99e-2}{mol}$
dan $m(\ce{Cu}) = 1.99\times 10^{-2} \times 63.5 = \qty{1.26}{g}$.
\end{solution}

\begin{solution}{exo:g11:redox:13}
\ce{Cr2O7^{2-} + 14H+ + 6e- -> 2Cr^{3+} + 7H2O} dan
$6\times$ (\ce{Fe^{2+} -> Fe^{3+} + e-}):
\ce{Cr2O7^{2-} + 14H+ + 6Fe^{2+} -> 2Cr^{3+} + 6Fe^{3+} + 7H2O}.
Enam \ce{Fe^{2+}} per \omterm{def:g9:ions:ion}{ion} dikromat: $6 \times 0.010 = \qty{0.060}{mol}$.
\end{solution}

\begin{solution}{exo:g11:redox:14}
\ce{2ClO- + 4H+ + 2e- -> Cl2 + 2H2O} dan \ce{2Cl- -> Cl2 + 2e-};
dijumlahkan lalu dibagi 2: \ce{ClO- + Cl- + 2H+ -> Cl2 + H2O}. Asam
menyediakan \omterm{def:g9:ions:ion}{ion-ion} \ce{H+}, dan reaksi itu melepaskan klorin, gas yang
beracun: itulah sebabnya ada peringatan agar kedua produk itu tidak
pernah \omterm{def:g2:mixing-and-dissolving:mix}{dicampur}.
\end{solution}

\begin{solution}{exo:g11:redox:15}
\omterm{def:g9:inside-the-atom:nucleus}{Elektron} tidak dapat bergerak bebas melalui larutan
(\cref{prop:g11:redox:no-free-electrons}): \omterm{def:g9:inside-the-atom:nucleus}{elektron} yang dilepaskan \omterm{def:g7:atoms-and-molecules:atom}{atom}
seng hanya dapat diambil oleh \omterm{def:g9:ions:ion}{ion} \ce{Cu^{2+}} yang menyentuh \omterm{def:g9:periodic-table-first-look:metal}{logam} itu.
Karena itu \omterm{def:g7:atoms-and-molecules:atom}{atom} tembaga terbentuk di permukaan keping itu, dan tetap di
sana.
\end{solution}

\begin{solution}{pb:g11:redox:1}
\textbf{1.} \ce{Cr2O7^{2-}}/\ce{Cr^{3+}} dan \ce{CH3COOH}/\ce{CH3CH2OH}.

\textbf{2.} Etanol dioksidasi: etanol adalah \omterm{def:g11:redox:oxidant}{reduktornya}.

\textbf{3.} \ce{Cr2O7^{2-} + 14H+ + 6e- <=> 2Cr^{3+} + 7H2O}.

\textbf{4.} Karbon sudah setara (dua di setiap ruas). Oksigen: dua di
kiri, satu di kanan, jadi \ce{H2O} di kanan. Hidrogen: empat di kiri,
$6 + 2 = 8$ di kanan, jadi \ce{4H+} di kiri. Muatan: $+4$ di kiri, 0 di
kanan, jadi empat \omterm{def:g9:inside-the-atom:nucleus}{elektron} di kiri:
\ce{CH3COOH + 4H+ + 4e- <=> CH3CH2OH + H2O}.

\textbf{5.} Dua belas \omterm{def:g9:inside-the-atom:nucleus}{elektron}: $2\times$ yang pertama, $3\times$ yang
kedua dibalik:
\ce{2Cr2O7^{2-} + 3CH3CH2OH + 16H+ -> 4Cr^{3+} + 3CH3COOH + 11H2O}.

\textbf{6.} \omterm{def:g9:ions:ion}{Ion-ion} \ce{Cr2O7^{2-}} yang jingga direduksi menjadi \omterm{def:g9:ions:ion}{ion}
\ce{Cr^{3+}} yang hijau di mana pun etanol mencapainya.

\textbf{7.} $M = 2\times 39.1 + 2\times 52.0 + 7\times 16.0 =
\qty{294.2}{g/mol}$.

\textbf{8.} $n = 2.0\times 10^{-3}/294.2 = \qty{6.8e-6}{mol}$.

\textbf{9.} Tiga \omterm{def:g7:atoms-and-molecules:molecule}{molekul} etanol per dua \omterm{def:g9:ions:ion}{ion} dikromat:
$n = \tfrac{3}{2}\times 6.8\times 10^{-6} = \qty{1.0e-5}{mol}$.

\textbf{10.} $M(\ce{C2H6O}) = \qty{46.0}{g/mol}$, jadi
$m = 1.02\times 10^{-5}\times 46.0 = \qty{4.7e-4}{g}$, sekitar
\qty{0.47}{mg}.

\textbf{11.} $0.25/0.47 \approx 0.53$: sekitar \qty{53}{\percent}
dikromatnya.

\textbf{12.} $0.53 \times 10 \approx \qty{5.3}{cm}$.

\textbf{13.} $\qty{0.25}{mg} \times 2100 = \qty{525}{mg}$ per liter
darah, sekitar \qty{0.53}{g/L}.

\textbf{14.} Napas itu lebih dulu bertemu butiran-butiran di lubang
masuk; dikromat di sana berlebih dan menghabiskan seluruh etanol,
sehingga reaksinya tuntas di bagian pertama tabung, yang menjadi hijau
sepenuhnya, sebelum etanol mencapai bagian yang lebih jauh.

\textbf{15.} Reaksi itu tuntas: \omterm{def:g9:ions:ion}{ion-ion} \ce{Cr^{3+}} yang hijau tidak
kembali menjadi dikromat, sehingga tabung yang sudah terpakai tidak dapat
mengukur lagi.

% ledger: ghs:K2Cr2O7 (H350 among its hazard statements)
\textbf{16.} Zat itu sangat beracun (GHS06) dan sangat berbahaya bagi
kesehatan (GHS08: pernyataan bahayanya mencakup ``dapat menyebabkan
kanker''), selain korosif dan beracun bagi kehidupan akuatik; alat yang
dipakai masyarakat umum tidak boleh mengandungnya.

\textbf{17.} \ce{O2 + 4H+ + 4e- -> 2H2O} dan
\ce{CH3CH2OH + H2O -> CH3COOH + 4H+ + 4e-}:
\ce{CH3CH2OH + O2 -> CH3COOH + H2O}.

\textbf{18.} Empat \omterm{def:g9:inside-the-atom:nucleus}{elektron}.

\textbf{19.} $n = 0.25\times 10^{-3}/46.0 = \qty{5.4e-6}{mol}$
etanol; \omterm{def:g9:inside-the-atom:nucleus}{elektronnya} $4n = \qty{2.2e-5}{mol}$, yaitu
$2.2\times 10^{-5}\times 6.02\times 10^{23} = \num{1.3e19}$ \omterm{def:g9:inside-the-atom:nucleus}{elektron}.

\textbf{20.} $Q = 1.3\times 10^{19}\times 1.60\times 10^{-19} \approx
\qty{2.1}{C}$. Setiap \omterm{def:g7:atoms-and-molecules:molecule}{molekul} etanol mengirimkan tepat empat \omterm{def:g9:inside-the-atom:nucleus}{elektron}
melalui kawat, sehingga muatannya sebanding dengan jumlah etanol:
mengukur muatan itu berarti mengukur etanolnya.
\end{solution}
